% Project report for the QSVM track, structured as a scientific paper.
%
% Build:  latexmk -pdf qsvm_project_report.tex
%
% Figures:
%   uv run python scripts/make_history_figures.py     -> figures_history/
%   uv run python scripts/make_report_figures.py      -> figures/
% Regenerate rather than editing them.
%
% Scope: this report presents results and the limitations OF THOSE RESULTS.
% Work that was never carried out is tracked in NOTES_open_items.md and is
% deliberately absent here.

\documentclass[11pt]{article}

\usepackage[a4paper,margin=1in]{geometry}
\usepackage{lmodern}
\usepackage[T1]{fontenc}
\usepackage[utf8]{inputenc}
\usepackage{amsmath,amssymb}
\usepackage{booktabs}
\usepackage{graphicx}
\usepackage{caption}
\usepackage{siunitx}
\usepackage[hidelinks]{hyperref}
\usepackage{microtype}
\usepackage{enumitem}

\captionsetup{font=small,labelfont=bf}
\sisetup{detect-weight=true,detect-inline-weight=math}
\setlist{itemsep=1pt,topsep=3pt}

\newcommand{\nc}{\ensuremath{n_{\mathrm{components}}}}
\newcommand{\best}[1]{\textbf{#1}}
\newcommand{\HFIG}[1]{figures_history/#1}
\newcommand{\WFIG}[1]{figures/#1}

\title{Does a fidelity-kernel quantum SVM outperform classical baselines\\
on malware family attribution?\\[0.3em]
\large A controlled comparison on two corpora}
\date{2 August 2026}

\begin{document}
\maketitle

\begin{abstract}
\noindent
\textbf{Background.} Quantum kernel methods are often reported as competitive
with classical machine learning on security data, but the two arms of such a
comparison frequently differ in the rows they are scored on, the input
dimensionality they receive, or the metric used to report the outcome.

\noindent
\textbf{Methods.} We compare a fidelity-kernel quantum support vector machine
(QSVM) against four tuned classical models on CIC-MalMem-2022 and EMBER 2018,
each with a binary and a 15-class task. Both frameworks consume one shared
feature pipeline whose final PCA dimensionality simultaneously sets the classical
feature count and the quantum qubit budget. The quantum runner writes its
stratified 1000-row subsample and its five outer folds to disk and the classical
runner replays them, so both are scored on identical rows. Differences are
assessed with a paired $t$-test over per-fold macro-F1.

\noindent
\textbf{Results.} On the binary task of CIC-MalMem every method exceeds
$0.9998$ macro-F1 on the full data, so that task cannot rank methods. On the
three configurations that do discriminate, classical models lead by $0.085$ to
$0.274$ macro-F1, and all 84 comparable pairs reach $p<0.05$. The QSVM falls to
the chance level on CIC 15-class. An RBF control built on the identical feature
matrix produces the same off-diagonal Gram spread on both corpora ($0.2212$
against $0.2184$), whereas the fidelity kernel matches RBF on EMBER
($0.92\times$) and collapses to $0.40\times$ on CIC.

\noindent
\textbf{Conclusions.} Under a controlled protocol the fidelity kernel does not
outperform tuned classical baselines on either corpus. Its collapse on CIC is a
property of the kernel rather than of the data, and kernel concentration is the
mechanism.
\end{abstract}

\section{Introduction}

Malware family attribution is a multiclass problem over feature vectors that are
high-dimensional, heavy-tailed and often near-degenerate. It has become a common
testbed for quantum kernel methods, in which a quantum feature map replaces a
classical kernel inside an otherwise standard support vector machine.

Claims of quantum competitiveness on such data are hard to evaluate, because
three things routinely differ between the quantum and classical arms of a
comparison: the rows each side is scored on, the dimensionality each side
receives, and the metric used to report the outcome. Any of the three can
manufacture or erase a difference.

This work asks one question under a protocol that removes those three degrees of
freedom: \emph{does a fidelity-kernel QSVM outperform tuned classical baselines
at malware family attribution?} The contributions are:

\begin{enumerate}
\item a comparison protocol in which a single parameter fixes the classical
feature count and the quantum qubit budget simultaneously, and in which the
quantum arm's subsample and cross-validation folds are persisted and replayed by
the classical arm (§\ref{sec:protocol});
\item results across two corpora, two tasks and three capacities, with paired
significance tests on every comparable pair (§\ref{sec:results});
\item evidence that one task widely used in this setting is saturated and
therefore cannot rank methods at all (§\ref{sec:saturation});
\item a diagnosis of the QSVM's collapse on one corpus that separates a property
of the data from a property of the kernel, using a classical control on the
identical feature matrix (§\ref{sec:kernel}).
\end{enumerate}

\section{Materials and methods}
\label{sec:methods}

\subsection{Datasets}

Two corpora, chosen so that no conclusion rests on a single feature
representation. Each supports a binary detection task and a 15-class family
attribution task.

\begin{description}[leftmargin=1.2em,style=nextline]
\item[CIC-MalMem-2022] Obfuscated-malware memory dumps described by 55 numeric
forensic features covering process, handle, callback and DLL counts. The binary
task uses all \num{58596} rows and is balanced between benign and malware. The
15-class task drops benign rows, since benign detection is what the binary task
measures, leaving \num{29298} malware rows over 15 families with a $1.7\times$
ratio between the largest and smallest.
\item[EMBER 2018] Static portable-executable features in 2381 dimensions. We work
on an \num{18014}-row extract of the published test set; the full release is a
single $200\,000 \times 2384$ block that cannot be read incrementally. The binary
task uses all \num{18014} rows and is $88.9\%$ malware. The 15-class task keeps
families with at least 500 samples, caps at the 15 largest, then downsamples each
to the smallest kept count, giving \num{14325} rows balanced at exactly 955 per
family.
\end{description}

The two differ in kind as well as in size. CIC describes runtime behaviour and
EMBER static file structure, and EMBER's families are exactly balanced where
CIC's are mildly skewed. That contrast is used in §\ref{sec:kernel}.

\subsection{Feature pipeline and the alignment parameter}

Both frameworks consume the output of one pipeline: a zero-variance filter, a
correlation filter at $0.95$, standardisation, then PCA to \nc{} components. The
pipeline is fit on training-fold rows only.

\nc{} is the single alignment point. It sets the number of features the classical
models receive \emph{and} the number of qubits the quantum circuit uses, so a
comparison at $\nc{}=6$ is six features against six qubits. Reporting a quantum
model on six qubits against a classical model on all 2381 raw features would not
be a comparison of methods.

\subsection{Quantum model}

The QSVM computes a fidelity kernel
$K(x_i,x_j) = |\langle \phi(x_j) | \phi(x_i) \rangle|^2$ by encoding $x_i$,
applying the inverse encoding of $x_j$, and reading the probability of the
all-zeros measurement outcome. The resulting Gram matrix is passed to a support
vector machine with a precomputed kernel. Two feature maps are compared: an angle
encoding with a ring of entangling gates, and an IQP (ZZ-interaction) encoding.
A bandwidth parameter scales inputs before rotation to counteract kernel
concentration and defaults to $n_{\mathrm{qubits}}^{-1/2}$.

All quantum results use a state-vector simulator.

\subsection{Comparison protocol}
\label{sec:protocol}

Building the training Gram requires one circuit evaluation per \emph{pair} of
samples, so its cost is $O(n^2)$ in the number of training rows.
Figure~\ref{fig:qcost} measures the per-pair component of that cost.

\begin{figure}[htbp]
\centering
\includegraphics[width=0.52\textwidth]{\HFIG{fig_quantum_cost.pdf}}
\caption{Cost of one circuit pair when building the training Gram, against qubit
count, taken from a single session so the worker configuration is constant.
Per-pair cost rises roughly tenfold from 1 to 8 qubits, and the number of pairs
grows as $n^2$, so the two effects compound.}
\label{fig:qcost}
\end{figure}

Full-dataset quantum runs are therefore infeasible: at $n=1000$ a single training
fold already requires $\approx\!\num{320000}$ circuit evaluations. Quantum results
use a stratified 1000-row subsample. A separate sizing exercise fitted the
scaling law and recommended $n=900$ at a 60\,s per-Gram budget, with a
class-balance floor of $n \ge 750$ for the 15-class task; the 1000-row subsample
sits just above that floor.

Because only the quantum arm is forced to subsample, the classical arm must be
scored on the same rows or the comparison is void. The quantum runner writes its
subsample indices and its five stratified outer folds to disk and the classical
runner reloads them, so neither side chooses its own split. Each fold provides
800 training and 200 held-out rows.

\subsection{Evaluation}

All results are held-out macro-F1, the unweighted mean of per-class F1, averaged
over the five outer folds. Hyperparameters are tuned by an Optuna inner loop
inside each training fold, followed by one refit on the full training fold.

Accuracy is not used. On EMBER binary at $\nc{}=1$ the QSVM attains $0.890$
accuracy in three of five folds while the benign class's F1 is exactly $0.000$ in
those same folds, because $88.9\%$ of rows are malware and the model labels
everything malware. Accuracy there reports near-success for a classifier that
never identifies the class of interest.

Differences between a quantum and a classical model are assessed with a paired
$t$-test over the five per-fold macro-F1 values.

\subsection{Selection of comparable runs}
\label{sec:selection}

Runs are grouped by cross-validation sweep rather than by configuration. The same
configuration was swept more than once over the course of the work, so a
parameter-based grouping returns 10, 15 or 20 records for a five-fold sweep and
averages distinct hyperparameter searches together; 57 parameter groups in the
experiment log have a size other than five. For the QSVM there is a second
reason: the encoding is a \emph{per-fold outcome} rather than a configuration
key, because the tuner selects it inside each fold, so a single sweep tuning over
both encodings can appear as four angle folds and one IQP fold.

Fold records are therefore grouped by parent run, retaining only sweeps that
completed with exactly five folds. Under that rule the log holds 125 nested
sweeps, 123 of them complete.

\section{Results}
\label{sec:results}

\subsection{The binary task on CIC-MalMem is saturated}
\label{sec:saturation}

\begin{figure}[htbp]
\centering
\includegraphics[width=\textwidth]{\HFIG{fig_w1_baseline.pdf}}
\caption{Four classical models on the full CIC-MalMem dataset without
dimensionality reduction. (a) On the binary task all four exceed $0.9998$
macro-F1 and the spread between best and worst is $10^{-4}$. (b) On a multiclass
formulation that still included benign as a class, the same models separate by
$0.19$ macro-F1.}
\label{fig:w1}
\end{figure}

Figure~\ref{fig:w1}(a) shows that the binary task on CIC-MalMem leaves no
headroom in which to measure a difference. Any ranking of methods on it measures
noise.

\begin{figure}[htbp]
\centering
\includegraphics[width=0.55\textwidth]{\HFIG{fig_cic_binary_ceiling.pdf}}
\caption{CIC-MalMem binary on the 1000-row subsample across the PCA sweep. Every
model, quantum and classical, lies between $0.977$ and $0.992$ macro-F1, and
every difference is inside the fold-to-fold error bars. Aggregated over all
recorded folds, since individual sweeps are not separable for this configuration
(§\ref{sec:limits}).}
\label{fig:ceiling}
\end{figure}

The saturation survives dimensionality reduction and subsampling
(Figure~\ref{fig:ceiling}). The QSVM does draw level with the classical models
here, at a gap of $0.006$ macro-F1. That is a property of the task rather than
evidence of parity between methods, and CIC binary is treated as a control
throughout.

\subsection{Classical models lead on every discriminating configuration}

\begin{table}[htbp]
\centering
\small
\caption{EMBER 2018, held-out macro-F1, mean over five outer folds. Best
classical model in bold.}
\label{tab:ember}
\begin{tabular}{@{}l S[table-format=1.4] S[table-format=1.4] S[table-format=1.4]
S[table-format=1.4] S[table-format=1.4] S[table-format=1.4]@{}}
\toprule
& {Random forest} & {XGBoost} & {LightGBM} & {SVM} & {QSVM angle} & {QSVM IQP} \\
\midrule
\multicolumn{7}{@{}l}{\emph{Binary}}\\
$\nc{}=1$ & \best{0.5582} & 0.5451 & 0.5093 & 0.5314 & 0.4532 & 0.4587 \\
$\nc{}=3$ & \best{0.6575} & 0.6465 & 0.6434 & 0.6306 & 0.4924 & 0.5084 \\
$\nc{}=6$ & \best{0.6835} & 0.6680 & 0.6593 & 0.6724 & 0.5138 & 0.5137 \\
\addlinespace
\multicolumn{7}{@{}l}{\emph{15-class}}\\
$\nc{}=1$ & \best{0.5562} & 0.5313 & 0.5410 & 0.5294 & 0.1688 & 0.1541 \\
$\nc{}=3$ & 0.7194 & 0.7089 & 0.7035 & \best{0.7389} & 0.3908 & 0.4797 \\
$\nc{}=6$ & 0.7903 & 0.7659 & 0.7709 & \best{0.7930} & 0.4232 & 0.5195 \\
\bottomrule
\end{tabular}
\end{table}

Table~\ref{tab:ember} gives the EMBER results. Which classical model leads
depends on the task: random forest on binary at every capacity, SVM on 15-class
at $\nc{} \in \{3,6\}$. Comparing the QSVM against SVM alone would understate the
classical advantage on binary by $0.011$ to $0.027$ macro-F1, which is why all
four baselines are reported.

\begin{figure}[htbp]
\centering
\includegraphics[width=\textwidth]{\WFIG{fig_capacity_scaling.pdf}}
\caption{Macro-F1 against the shared feature and qubit budget. Solid lines are
classical, dashed lines the two QSVM encodings. Error bars are one standard
deviation over the five outer folds, and the dotted line marks the 15-class
chance level.}
\label{fig:capacity}
\end{figure}

Figure~\ref{fig:capacity} shows behaviour across capacity. Both frameworks
improve monotonically but at different rates: on EMBER binary the QSVM gains
$0.055$ macro-F1 from one to six components where the random forest gains
$0.125$, so the gap widens with capacity rather than closing. On CIC 15-class the
quantum curves sit on the chance level at every capacity tested.

\begin{figure}[htbp]
\centering
\includegraphics[width=0.62\textwidth]{\HFIG{fig_project_arc.pdf}}
\caption{Best classical against best QSVM in each corpus and task, over
$\nc{} \in \{1,3,6\}$. Hatched bars are aggregated from raw folds because
individual sweeps are not separable for that configuration.}
\label{fig:arc}
\end{figure}

Summarising (Figure~\ref{fig:arc}): the classical lead is $0.085$ on CIC
15-class, $0.170$ on EMBER binary and $0.274$ on EMBER 15-class. On the saturated
control the two are $0.006$ apart.

\subsection{The differences are statistically significant}

\begin{figure}[htbp]
\centering
\includegraphics[width=0.56\textwidth]{\WFIG{fig_significance.pdf}}
\caption{Effect size against paired $t$-test $p$-value for all 84 comparable
quantum--classical pairs. Every point lies left of $p=0.05$ and above zero.}
\label{fig:sig}
\end{figure}

All 84 comparable pairs favour the classical model and all 84 reach $p<0.05$ on
the paired $t$-test. The largest $p$-value across every pair is $0.0241$ and the
smallest is $4 \times 10^{-6}$.

A Wilcoxon signed-rank test returns exactly $0.0625$ for all 84 pairs. With five
paired samples that is the smallest attainable two-sided $p$-value, so the test
cannot reach $0.05$ at this fold count regardless of effect size; all 84 reaching
the floor indicates that every pair achieved the most extreme rank configuration
available.

\subsection{The encoding advantage is corpus-dependent}

\begin{figure}[htbp]
\centering
\includegraphics[width=\textwidth]{\HFIG{fig_encoding_verdict.pdf}}
\caption{IQP minus angle encoding on the 15-class task. The grey band is the
fold-to-fold spread of the encodings themselves, so a difference inside that band
is within noise.}
\label{fig:enc}
\end{figure}

IQP is never worse than angle, but the margin is resolvable on only one corpus
(Figure~\ref{fig:enc}). On CIC the differences are $+0.0009$, $+0.0122$,
$+0.0113$ and $+0.0046$ at $\nc{}=1,3,6,8$, each smaller than either encoding's
own standard deviation. On EMBER the differences at $\nc{}=3$ and $6$ are
$+0.0889$ and $+0.0963$, well outside the corresponding spread. Extending CIC
15-class to 8 qubits did not lift performance above chance, which excludes qubit
count as the limiting factor there.

\subsection{The subsample preserves its population}

\begin{figure}[htbp]
\centering
\includegraphics[width=0.55\textwidth]{\WFIG{fig_representativeness.pdf}}
\caption{Per-feature two-sample Kolmogorov--Smirnov rejections at $\alpha=0.05$,
1000-row subsample against full population, beside the count chance alone
produces.}
\label{fig:repr}
\end{figure}

Since every quantum result is computed on 1000 rows, the subsample's fidelity to
its population is a precondition for the rest. A rejection count is interpretable
only against its null expectation, because a faithful subsample still rejects
about $\alpha \cdot n_{\mathrm{features}}$ features by chance. CIC rejects 0 of 55
features where $\approx 2.75$ were expected, and EMBER 1 of 2381 where
$\approx 119$ were expected (Figure~\ref{fig:repr}). Maximum class-proportion
drift is below $0.001$ on both.

\subsection{Kernel diagnostics}
\label{sec:kernel}

The QSVM reaches $0.056$--$0.079$ on CIC 15-class against a chance level of
$\approx 0.067$, while the same method reaches $0.42$--$0.52$ on EMBER 15-class.
Class imbalance is excluded, since EMBER's families are exactly balanced and
CIC's only $1.7\times$ skewed, and qubit count is excluded by the preceding
section. Two further measurements separate a property of the data from a property
of the kernel.

\paragraph{Feature geometry.}

\begin{figure}[htbp]
\centering
\includegraphics[width=\textwidth]{\WFIG{fig_geometry.pdf}}
\caption{Class separability of the shared PCA projection, averaged over the five
outer folds. Both metrics are scale-invariant and therefore comparable across
corpora of very different dimensionality.}
\label{fig:geom}
\end{figure}

On the identical projected features, EMBER's families are $\approx 8\times$ more
separable by one-vs-rest Fisher ratio, and up to $50\times$ on the worst-case
family. CIC's silhouette coefficient is negative at every \nc{}: the average point
lies closer to another family's centroid than to its own. CIC's PCA nonetheless
explains substantially more variance ($33$--$81\%$ against $8$--$30\%$), so its
variance is concentrated in directions that do not discriminate family.

This accounts for the poor absolute performance of every method on CIC 15-class.
It does not account for the QSVM-specific deficit, since a classical SVM extracts
$0.108$--$0.165$ macro-F1 from these same features where the QSVM extracts
$0.056$--$0.079$.

\paragraph{Kernel geometry.}

\begin{figure}[htbp]
\centering
\includegraphics[width=\textwidth]{\WFIG{fig_kernel_diagnostics.pdf}}
\caption{First outer training fold, $\nc{}=6$, IQP encoding, 800 rows, with an
RBF control built on the identical feature matrix. (a) Off-diagonal Gram spread.
(b) Kernel-target alignment relative to its floor: a constant kernel scores
$1/\sqrt{15}\approx 0.258$ on 15 balanced classes, so only the excess is
informative.}
\label{fig:kernel}
\end{figure}

The RBF control produces near-identical off-diagonal spread on both corpora
($0.2212$ against $0.2184$), so CIC's features do not cause a classical kernel to
concentrate. The fidelity kernel tracks RBF on EMBER at $0.92\times$ and collapses
to $0.40\times$ on CIC: every off-diagonal similarity is compressed into a band
$2.5\times$ narrower from the same inputs, driving the Gram toward a scaled
identity.

On alignment, the CIC fidelity kernel lands $0.000069$ \emph{below} the
constant-kernel floor, making it indistinguishable from a matrix of ones, and the
RBF control is also below it. On EMBER both clear the floor by $\approx 0.03$.
Alignment therefore separates the two corpora, and concentration separates the
two kernels.

Across the full sweep the CIC fidelity Gram is $1.9$--$2.6\times$ less spread than
EMBER's at every capacity and for both encodings, and concentration worsens with
qubit count on both, which is the trend the bandwidth parameter exists to
counteract.

\section{Discussion}

\subsection{Interpretation}

The central result is negative and consistent. Under a protocol that equalises
rows, folds, dimensionality and metric, the fidelity-kernel QSVM does not
outperform tuned classical baselines on either corpus or either task, and the
deficit is significant on every comparable pair.

The one configuration in which the two are level is the saturated CIC binary
task, where all methods sit within $0.015$ macro-F1 of one another. Reports of
quantum parity on binary malware detection should be read with that ceiling in
mind: a task on which four tuned classical models differ by $10^{-4}$ cannot
distinguish a quantum kernel from a classical one either.

The more transferable finding concerns \emph{why} the QSVM fails on one corpus
and not the other, and the RBF control isolates it. Since a classical kernel on
the identical feature matrix does not concentrate, the collapse is not a
consequence of CIC's features being hard; it is a consequence of the fidelity
kernel's behaviour on those features. A Gram matrix approaching a scaled identity
presents the support vector machine with a geometry in which all pairs are
equally similar, which is what a macro-F1 pinned to chance looks like internally.
This predicts that the method will struggle wherever the encoded states become
mutually near-orthogonal, independently of the corpus.

\subsection{Limitations}
\label{sec:limits}

\begin{description}[leftmargin=1.2em,style=nextline]
\item[Scale] Every quantum result derives from 1000 rows, forced by the quadratic
kernel cost, against source corpora of \num{29298} and \num{14325} rows for the
15-class tasks. The subsample is verified faithful, but the reported standard
deviations are over five folds of 200 held-out rows.
\item[Bandwidth is untested on CIC] Concentration is the identified mechanism and
bandwidth is the parameter that acts on it, but no bandwidth sweep was run on
CIC. §\ref{sec:kernel} identifies a mechanism; it does not demonstrate a remedy.
\item[Alignment is a weak predictor] Alignment reports that neither kernel carries
label information on CIC, yet a classical SVM reaches $0.16$ there. Alignment is a
global, unweighted average over all pairs, and a support vector machine requires
no such global property, only support vectors, a regularisation constant and
class weights. It ranks corpus difficulty usefully and predicts attainable SVM
performance poorly. The concentration measurement does not share this weakness,
since it describes the matrix the SVM actually receives.
\item[One quantum method] The comparison is a fidelity-kernel QSVM against
classical baselines, not quantum methods in general.
\item[Simulator only] All quantum results use a state-vector simulator, with no
hardware execution and no noise model. Concentration behaviour on hardware, where
sampling noise adds its own floor, is untested.
\item[One configuration is not separable] Every CIC binary run, classical and
quantum alike, predates sweep-level tracking in the experiment log and carries no
recoverable sweep boundary; 45 fold records per model resolve into two distinct
sweeps per \nc{} with no identifier separating them. Those records are valid
measurements and are used descriptively in Figures~\ref{fig:ceiling}
and~\ref{fig:arc}, but no grouping was reconstructed, since the time gaps between
runs are not a recorded boundary. The limitation is symmetric across frameworks.
\item[No sample-level test between frameworks] McNemar's test requires per-sample
predictions from both models on a shared test set. QSVM predictions were not
persisted during the reported sweeps, so cross-framework evidence rests on five
paired fold-level scores rather than on paired sample-level outcomes.
\item[Fold count] Five outer folds support the paired $t$-test and are too few for
Wilcoxon, whose minimum attainable $p$-value at $n=5$ is $0.0625$. Six folds would
place $0.03125$ within reach.
\end{description}

\subsection{Future work}

Two steps follow directly from §\ref{sec:kernel}. First, sweep bandwidth on CIC:
the mechanism is identified and the corresponding parameter has never been tuned
there beyond its default. Second, move to six outer folds, which makes a
non-parametric test usable and removes the ceiling noted above.

\section{Conclusion}

Under a protocol that holds rows, folds, dimensionality and metric fixed across
both arms, a fidelity-kernel QSVM does not outperform tuned classical baselines at
malware family attribution on either of two corpora. Classical models lead on
every configuration where the task can discriminate at all, by $0.085$ to $0.274$
macro-F1, with $p<0.05$ on all 84 comparable pairs. The single level result occurs
on a task where all methods are saturated and no method can be ranked.

The QSVM's collapse to chance on one corpus is attributable to kernel
concentration rather than to feature difficulty: a classical RBF kernel on the
identical feature matrix retains its off-diagonal spread while the fidelity kernel
falls to $0.40\times$ of it. That is a statement about the method rather than
about the data, and it identifies both the failure mode and the parameter that
acts on it.

\end{document}
